\scanpage{30}

On a

\[
N(\Lambda\sigma(\Lambda)^{-1})\simeq N(\Lambda)(N(\sigma(\unsure{A|\Lambda})))^{-1}\simeq\underline O,
\]

isom can dû au fait que

\[
N(\sigma(\unsure{A|\Lambda}))\simeq N(\Lambda)\qquad\text{canoniquement.}
\]

Donc on s'attend à trouver un $A$-Module loc libre, avec restriction du groupe structural de $W(A)^*$ à $U(A)$, \inserted{schéma en} groupes des sections locales de $A$ de norme $1$.

Le $\mathcal A$-Module qui se propose à la vue est évidemment $\mathcal M/\mathcal A$, lequel est bien libre de rang $1$ sur $\mathcal A$, si $\mathcal A$ est un sous-fibré vect. de rang $2$ \deleted{de $\mathcal M$} (cas ``str. fidèle''). \deleted{Mais}

Dans le cas où $\mathcal M=\underline{\operatorname{End}}(E)$, il faudrait voir qu'on a un ``iso. à scalaire près'' canonique

\[
E\xrightarrow{\sim}\underline{\operatorname{End}}(E)/\mathcal A
\]

(iso. de $\mathcal A$-modules, mais le ``scalaire près'' réfère à une section de $\underline O_X^*$ \ldots), \deleted{ou encore un iso de schémas}

\[
\deleted{{\unsure{P|\mathbb P}}^{\unsure{\vee|*}}(E)\simeq P\bigl(\underline{\operatorname{End}}(E)^{\unsure{*|+}}/\mathcal A\bigr)}
\]

% left margin, vertical strokes
\margin{$\begin{gathered}|\\|\\|\\\unsure{!|\vert}\\)\end{gathered}$}

% slanted
\margin{\deleted{où $\underline{\operatorname{End}}^{+}$ désigne le sous-Module des end. de \uline{trace nulle}}}

\deleted{ou encore (ou mieux}

\[
\begin{aligned}
&\deleted{P(E)\simeq\unsure{\mathcal Q}\{\text{schéma des droites vect. }D\text{ dans }\underline{\operatorname{End}}(E)^+\text{ telles que }N\mid D=0\}},\\
&\deleted{D\longmapsto D^{\perp}\otimes D}
\end{aligned}
\]
% archive sheet number: 29
